% =====================================================================
% 01_environment.tex — the simplified environment (Definition 1) (WP-A)
% =====================================================================
\section{The simplified environment}
\label{sec:environment}

The analytical results of KL's Section~II are derived in a \emph{simplified
environment}: a set of parametric restrictions on the full model of Section~I
that admits closed-form first-order results while preserving the two essential
ingredients --- a time-varying convenience yield on safe dollar bonds and
cross-country heterogeneity in risk-bearing capacity. We state these
restrictions exactly as KL's Definition~1 (their p.~1660), and after each item
note which step of the analysis it buys, following KL's own discussion on
p.~1660.

\begin{definition}[Simplified environment; KL Definition~1]
\label{def:simplified}
The simplified environment features:
\begin{enumerate}[label=(\roman*),leftmargin=2.0em,itemsep=3pt]
  \item \textbf{Flexible wages, or wages set one period in advance.} In place of
  the Rotemberg (1982) wage-adjustment cost of the full model (their eq.~(5),
  scale $\wadj$), wages are either fully flexible or predetermined by one
  period. \\[2pt]
  \emph{What it buys:} this replaces the forward-looking wage Phillips curve
  (their app.\ eq.~(52)) by a static or one-period-lagged wage equation,
  simplifying the dynamics so that the impulse responses can be solved in closed
  form. The two cases (flexible vs.\ one-period-in-advance) correspond exactly
  to the two bullet groups in Propositions~1--2.
  \item \textbf{Fixed global capital stock} $\bigl(\kadj\!\to\!\infty,\
  \dep\!\to\!0\bigr)$. Infinite capital-adjustment costs freeze the global
  capital stock $\bar k_t$, and zero depreciation removes investment dynamics. \\[2pt]
  \emph{What it buys:} together with (i), this eliminates the endogenous
  capital-accumulation state and the investment Euler equation as a source of
  dynamics, so the only ``real'' margins left are the international allocation of
  a fixed capital stock and the portfolio. It also makes capital a pure
  store-of-value claim whose quantity is fixed.
  \item \textbf{Unitary IES} $(\ies=1)$, \textbf{complete home bias}
  $\bigl(\hbias\!\to\!\zs/(1+\zs)\bigr)$, and \textbf{infinite Frisch
  elasticity} $(\frisch\!\to\!0)$. \\[2pt]
  \emph{What it buys:} $\ies=1$ collapses the Epstein--Zin time aggregator to the
  Cobb--Douglas / log case, so that intertemporal substitution drops out of the
  consumption Euler equation's drift and the value-function recursion
  log-linearizes cleanly. Complete home bias
  $\hbias\!\to\!\zs/(1+\zs)$ sends the CES weight on the foreign good in
  consumption to zero (each country consumes only its own good in the
  deterministic steady state), which decouples the steady-state real exchange
  rate from the consumption shares and is what makes the role of home bias
  \emph{transparent} in the final results. Infinite Frisch ($\frisch\!\to\!0$ in
  the parameterization of their eq.~(3)) makes the labor-disutility multiplier
  $\Phi(\ell)$ linear in the relevant sense, simplifying the supply block.
  \item \textbf{No disaster risk} $(p=0,\ \sigma^p=0)$, \textbf{constant relative
  productivity} $(\sigma^F=0)$, and \textbf{transitory safety}
  $(\rho^\omega=0)$. \\[2pt]
  \emph{What it buys:} switching off disasters and Foreign-productivity shocks
  leaves \emph{two} aggregate shocks --- global productivity $\epsilon^z$ and the
  safety shock $\epsilon^\omega$ --- which is exactly the count needed for the
  ``locally complete'' three-asset structure used in the portfolio analysis
  (Proposition~4). Transitory safety ($\rho^\omega=0$) means a safety shock dies
  out after one period, which is what licenses the period-$1$/period-$\ge 2$
  decomposition of the impulse responses (their app.\ B.1).
  \item \textbf{Identical per-capita wealth across countries in the
  deterministic steady state.} \\[2pt]
  \emph{What it buys:} this pins down a symmetric deterministic steady state
  around which the perturbation is taken, ensuring the steady state is
  well-defined and that the linearized system has the symmetric structure used
  throughout the proofs.
  \item \textbf{Identical discount factors} $(\disc=\discF)$. \\[2pt]
  \emph{What it buys:} with $\disc=\discF$ the two countries share the same
  steady-state interest rate and (given (v)) the same steady-state wealth share
  $\theta=1/(1+\zs)$, removing a level asymmetry that would otherwise complicate
  the steady state. (In the calibrated model KL allow $\disc\neq\discF$ to match
  the level of net foreign assets; here it is shut off.)
\end{enumerate}
\end{definition}

\begin{remark}[General home bias in the proofs; KL app.\ B.1]
\label{rem:general-sigma}
A crucial nuance for what follows: although Definition~\ref{def:simplified}
imposes \emph{complete} home bias $\hbias\!\to\!\zs/(1+\zs)$, KL carry out the
impulse-response and portfolio proofs of Appendix~B for a \emph{general}
home-bias parameter $\hbias$, and only specialize to complete home bias
$\hbias\!\to\!\zs/(1+\zs)$ at the very end, ``so that the role of home bias is
clear'' (their app.\ B.1, opening paragraph). Accordingly, the re-scaled
equilibrium system (Section~\ref{sec:setup} below) and the period-$1$
log-linearization (introduced in the impulse-response engine, WP-B) retain
$\hbias$ as a free parameter; the limit $\hbias\!\to\!\zs/(1+\zs)$ is taken only
where a proposition's stated form requires it. We preserve this convention
throughout this reference: \emph{$\hbias$ is general until explicitly
specialized}.
\end{remark}

\begin{remark}[Perturbation, not the global solution]
The simplified environment is studied with a perturbation (first-order for the
impulse responses, second-order for portfolios) around the deterministic steady
state guaranteed by (v)--(vi). KL emphasize that \emph{none} of the assumptions
in Definition~\ref{def:simplified} is imposed in the quantitative sections~III--V,
which instead solve the full model globally (their app.\ A.7). This reference
concerns only the analytical (simplified-environment) results.
\end{remark}
